\documentclass[11pt]{article}

% ── Packages ──────────────────────────────────────────────────────────────────
\usepackage[T1]{fontenc}
\usepackage[utf8]{inputenc}
\usepackage{lmodern}
\usepackage[margin=1in]{geometry}
\usepackage{microtype}
\usepackage{amsmath,amssymb}
\usepackage{booktabs}
\usepackage{multirow}
\usepackage{graphicx}
\usepackage{xcolor}
\usepackage{listings}
\usepackage{enumitem}
\usepackage[hidelinks,
            colorlinks=true,
            linkcolor=blue!60!black,
            citecolor=blue!60!black,
            urlcolor=blue!60!black]{hyperref}
\usepackage{fancyhdr}
\usepackage{titlesec}
\usepackage{abstract}

% ── Typography ────────────────────────────────────────────────────────────────
\titleformat{\section}{\large\bfseries}{\thesection.}{0.5em}{}
\titleformat{\subsection}{\normalsize\bfseries}{\thesubsection.}{0.5em}{}
\setlength{\parskip}{3pt}
\setlength{\parindent}{1em}

% ── Code listings ─────────────────────────────────────────────────────────────
\lstset{
  basicstyle=\ttfamily\footnotesize,
  breaklines=true,
  frame=single,
  framerule=0.4pt,
  columns=flexible,
  keepspaces=true,
  backgroundcolor=\color{gray!8},
}

% ── Header / Footer ───────────────────────────────────────────────────────────
\pagestyle{fancy}
\fancyhf{}
\renewcommand{\headrulewidth}{0.4pt}
\fancyhead[L]{\footnotesize ANE-LM Batch Prefill Bench}
\fancyhead[R]{\footnotesize AtomGradient, 2026}
\fancyfoot[C]{\thepage}

% ── Title ─────────────────────────────────────────────────────────────────────
\title{\textbf{ANE Batch Prefill for On-Device Parallel LLM Inference:\\
       Enabling Concurrent ANE Prefill and GPU Decode on Apple Silicon}}

\author{%
  AtomGradient\\
  \small\href{https://github.com/AtomGradient/hybird-batch-prefill-on-ane}{github.com/AtomGradient/hybird-batch-prefill-on-ane}
}

\date{March 2026}

% ══════════════════════════════════════════════════════════════════════════════
\begin{document}
\maketitle

% ── Abstract ──────────────────────────────────────────────────────────────────
\begin{abstract}
Our prior work~\cite{atomgradient2026hybrid} demonstrated that Apple
Silicon's Neural Engine (ANE) can offload LLM prefill from the GPU, reducing
GPU power draw by 282$\times$ and freeing thermal budget for sustained decode.
However, the private \texttt{AppleNeuralEngine.framework} API used in that
study dispatches only one token per call, yielding a constant latency of
$\sim$42\,ms/token --- impractical for interactive inference.

This companion paper addresses that bottleneck by introducing \textbf{ANE
batch prefill}: processing up to 32 tokens per ANE dispatch
(\texttt{ANE\_SPATIAL=32}) via fused matrix--vector kernels compiled at
model-load time. We benchmark three inference modes --- pure ANE, pure MLX
(GPU), and hybrid (ANE batch prefill + MLX GPU decode) --- on Qwen3.5-0.8B
and Qwen3.5-2B across short (13 tokens), medium (28 tokens), and long (74
tokens) prompts on an M2 Ultra.

Key results:
(i)~ANE batch prefill reaches 268\,tok/s on 0.8B (vs.\ 23.7\,tok/s for
sequential dispatch --- an \textbf{11.3$\times$ speedup});
(ii)~ANE batch prefill \emph{exceeds} GPU prefill throughput on short
prompts (149 vs.\ 138\,tok/s for 0.8B at 13 tokens);
(iii)~hybrid decode throughput reaches 39--54\,tok/s on 0.8B (vs.\
$\sim$18\,tok/s for pure ANE decode);
(iv)~state transfer overhead (save + read + cache construction) is
$<$30\,ms total, negligible relative to inference time.
These results validate the feasibility of \textbf{concurrent ANE prefill +
GPU decode pipelines} for multi-request serving on edge devices.
\end{abstract}

% ══════════════════════════════════════════════════════════════════════════════
\section{Introduction}
\label{sec:intro}

On-device LLM inference faces a fundamental tension: the prefill phase is
compute-bound (matrix multiplications over the full prompt), while the
decode phase is memory-bandwidth-bound (sequential token generation).
Apple Silicon provides three compute units --- CPU, GPU, and Neural Engine
(ANE) --- sharing a unified memory bus, offering an opportunity to
\emph{disaggregate} these phases across hardware.

Our prior work~\cite{atomgradient2026hybrid} explored four inference
pipelines on the same hardware, concluding that:
\begin{itemize}[nosep]
  \item CoreML batched ANE prefill reaches GPU parity at $\sim$410 tokens,
  \item Private API sequential dispatch has a constant $\sim$42\,ms/token
        latency (the fundamental bottleneck),
  \item ANE prefill reduces GPU power from 62.05\,W to 0.22\,W.
\end{itemize}

The critical limitation was \textbf{sequential dispatch}: the private
\texttt{AppleNeuralEngine.framework} API processes only one token per
call, making prefill throughput scale as $O(N)$ regardless of ANE
hardware capability. CoreML's batched dispatch solves this but requires
a heavy framework layer and model compilation step.

This paper demonstrates that the private API \emph{can} achieve batched
prefill by compiling fused ANE kernels that process up to 32 tokens
simultaneously (\texttt{ANE\_SPATIAL=32}). The key contributions are:

\begin{enumerate}[nosep]
  \item \textbf{ANE batch prefill implementation}: Fused matrix--vector
        kernels compiled at model-load time, processing 32 tokens per
        dispatch while CPU handles per-token state updates (RMSNorm, RoPE,
        SSM recurrence, GQA attention).
  \item \textbf{Hybrid pipeline validation}: ANE batch prefill + MLX GPU
        decode via binary state serialization (64-byte header + per-layer
        state), demonstrating $<$30\,ms transfer overhead.
  \item \textbf{Throughput characterization}: Benchmarks across two model
        sizes (0.8B, 2B) and three prompt lengths, showing batch prefill
        achieves 11.3$\times$ speedup over sequential dispatch.
  \item \textbf{Concurrent pipeline feasibility}: Since ANE batch prefill
        leaves the GPU idle (0.22\,W measured in~\cite{atomgradient2026hybrid}),
        GPU decode of a previous request can execute \emph{simultaneously},
        enabling pipelined multi-request serving.
\end{enumerate}

% ══════════════════════════════════════════════════════════════════════════════
\section{Background}
\label{sec:background}

\subsection{Hardware Platform}

All experiments were conducted on a Mac Studio with an Apple M2 Ultra
(Table~\ref{tab:hardware}).

\begin{table}[h]
\centering
\caption{Test platform specifications.}
\label{tab:hardware}
\begin{tabular}{ll}
\toprule
Spec & Value \\
\midrule
Chip           & Apple M2 Ultra \\
CPU cores      & 24 (16P + 8E) \\
GPU cores      & 76 \\
ANE cores      & 32 \\
ANE TOPS       & 31.6 \\
Memory         & 192\,GB unified \\
Bandwidth      & 800\,GB/s \\
Metal support  & Metal 4 \\
\bottomrule
\end{tabular}
\end{table}

\subsection{Model Architecture: Qwen3.5 Hybrid}

Qwen3.5 uses a hybrid architecture combining \textbf{DeltaNet} linear
attention layers with periodic \textbf{full attention} layers (every
\texttt{full\_attention\_interval} layers). This creates two distinct
state types:

\begin{itemize}[nosep]
  \item \textbf{DeltaNet layers}: SSM state $[H, K, V]$ + circular
        convolution buffer $[C, \text{kernel}-1]$
  \item \textbf{Full attention layers}: KV cache
        $[\text{capacity}, \text{heads}, \text{dim}]$ with circular
        buffer (capacity = 2{,}048)
\end{itemize}

The 0.8B model has 24 layers (18 DeltaNet + 6 FullAttn); the 2B model
has 36 layers (27 DeltaNet + 9 FullAttn).

\subsection{Sequential vs.\ Batched ANE Dispatch}

Our prior work~\cite{atomgradient2026hybrid} identified two dispatch
modes for ANE inference:

\textbf{Sequential dispatch} (prior work, Pipelines 3--4): Each token
requires a separate ANE model descriptor evaluation
via \texttt{IOSurface} shared memory. Dispatch latency is
$\sim$42\,ms/token regardless of prompt length, yielding
TTFT~$\propto N$.

\textbf{Batched dispatch} (this work): ANE kernels are compiled to process
\texttt{ANE\_SPATIAL=32} tokens per evaluation. The ANE performs the
matrix multiplications for all 32 tokens in parallel, while the CPU
sequentially updates per-token state (RMSNorm, RoPE, SSM recurrence,
attention). This reduces ANE dispatch overhead by up to 32$\times$.

% ══════════════════════════════════════════════════════════════════════════════
\section{Methods}
\label{sec:methods}

\subsection{ANE Batch Prefill}

The inference pipeline splits each transformer layer between ANE and CPU:

\begin{enumerate}[nosep]
  \item \textbf{ANE}: Executes all linear projections as fused
        matrix--vector operations (\texttt{first\_proj}, \texttt{o\_proj},
        \texttt{fused\_ffn}, \texttt{lm\_head}). Kernels are compiled to
        accept spatial dimension = 32, enabling batch processing.
  \item \textbf{CPU}: Executes non-linear operations sequentially per
        token: RMSNorm (via \texttt{vDSP}), RoPE, SSM recurrence
        (DeltaNet state update), Conv1d, GQA attention with KV cache, and
        softmax (all via Accelerate framework CBLAS/vDSP).
\end{enumerate}

Prompt tokens are processed in batches of up to 32. For a 74-token prompt,
this requires $\lceil 74/32 \rceil = 3$ ANE dispatches (vs.\ 74 sequential
dispatches in the prior approach). ANE kernels are compiled once at model
load and cached persistently to \texttt{\{model\_dir\}/ane\_weights/}.

\subsection{State Serialization}

To bridge ANE prefill output to MLX GPU decode, we serialize the full
model state to a binary file:

\begin{itemize}[nosep]
  \item \textbf{Header} (64 bytes): Magic (\texttt{ANELMS\textbackslash0\textbackslash0}),
        version, layer count, prompt length, logits dimension.
  \item \textbf{Per-layer state}: DeltaNet layers store
        \texttt{conv\_pos} + SSM state + conv state; full attention layers
        store \texttt{kv\_len} + \texttt{kv\_start} + K cache + V cache.
  \item \textbf{Footer}: Final logits vector for next-token sampling.
\end{itemize}

The MLX decode script reads this binary state, constructs
\texttt{ArraysCache} (DeltaNet) or \texttt{KVCache} (FullAttn)
objects, and continues autoregressive generation on the GPU.

\subsection{Three Inference Modes}

We benchmark three modes:
\begin{enumerate}[nosep]
  \item \textbf{Pure ANE}: Both prefill and decode on ANE+CPU (batch
        prefill for prompt, sequential decode for generation).
  \item \textbf{Pure MLX}: Both prefill and decode on GPU via MLX
        framework (baseline).
  \item \textbf{Hybrid}: ANE batch prefill $\rightarrow$ state
        serialization $\rightarrow$ MLX GPU decode. The state transfer
        includes save (ANE side), read (MLX side), and cache construction.
\end{enumerate}

% ══════════════════════════════════════════════════════════════════════════════
\section{Experiments}
\label{sec:experiments}

\subsection{Setup}

We tested two Qwen3.5 models:
\begin{itemize}[nosep]
  \item \textbf{Qwen3.5-0.8B} (FP16, 24 layers, $d=1{,}024$)
  \item \textbf{Qwen3.5-2B} (BF16, 36 layers, $d=1{,}536$)
\end{itemize}

Three prompt lengths were used:
\begin{itemize}[nosep]
  \item \textbf{Short}: ``Hello'' (13 tokens with chat template)
  \item \textbf{Medium}: ``Explain the theory of relativity\ldots'' (28 tokens)
  \item \textbf{Long}: ``Write a detailed technical comparison\ldots'' (74 tokens)
\end{itemize}

Each configuration was run 10 times; we report averages. Decode was capped
at 20 tokens with temperature = 0 (greedy sampling).

\subsection{Throughput Results: Qwen3.5-0.8B}

\begin{table}[h]
\centering
\caption{Throughput comparison for Qwen3.5-0.8B FP16 on M2 Ultra.
  All values averaged over 10 runs.}
\label{tab:results-08b}
\begin{tabular}{llrrr}
\toprule
Prompt & Mode & Prefill (tok/s) & Decode (tok/s) & Total (ms) \\
\midrule
\multirow{3}{*}{Short (13 tok)}
  & ANE      & 148.7 & 17.6 & 2{,}931 \\
  & MLX      & 138.2 & 61.2 & 3{,}679 \\
  & Hybrid   & 148.3 & 39.4 & 6{,}050 \\
\midrule
\multirow{3}{*}{Medium (28 tok)}
  & ANE      & 237.5 & 16.4 & 3{,}645 \\
  & MLX      & 290.0 & 64.5 & 3{,}836 \\
  & Hybrid   & 237.9 & 53.9 & 6{,}280 \\
\midrule
\multirow{3}{*}{Long (74 tok)}
  & ANE      & 268.0 & 17.1 & 3{,}774 \\
  & MLX      & 736.2 & 64.1 & 3{,}846 \\
  & Hybrid   & 271.1 & 54.0 & 6{,}335 \\
\bottomrule
\end{tabular}
\end{table}

\subsection{Throughput Results: Qwen3.5-2B}

\begin{table}[h]
\centering
\caption{Throughput comparison for Qwen3.5-2B BF16 on M2 Ultra.
  All values averaged over 10 runs.}
\label{tab:results-2b}
\begin{tabular}{llrrr}
\toprule
Prompt & Mode & Prefill (tok/s) & Decode (tok/s) & Total (ms) \\
\midrule
\multirow{3}{*}{Short (13 tok)}
  & ANE      & 91.7  & 12.4 & 5{,}271 \\
  & MLX      & 179.6 & 90.8 & 4{,}076 \\
  & Hybrid   & 92.8  & 22.6 & 8{,}878 \\
\midrule
\multirow{3}{*}{Medium (28 tok)}
  & ANE      & 160.4 & 12.5 & 6{,}189 \\
  & MLX      & 358.0 & 94.1 & 4{,}212 \\
  & Hybrid   & 159.3 & 29.0 & 9{,}281 \\
\midrule
\multirow{3}{*}{Long (74 tok)}
  & ANE      & 172.7 & 12.7 & 6{,}455 \\
  & MLX      & 828.6 & 92.4 & 4{,}274 \\
  & Hybrid   & 173.1 & 29.2 & 9{,}518 \\
\bottomrule
\end{tabular}
\end{table}

\subsection{State Transfer Overhead}

\begin{table}[h]
\centering
\caption{State transfer overhead for hybrid pipeline (averaged over all
  prompt lengths and runs).}
\label{tab:overhead}
\begin{tabular}{lrrr}
\toprule
Model & Save (ms) & Read (ms) & Cache (ms) \\
\midrule
0.8B & 12 & 9 & 6 \\
2B   & 11 & 9 & 6 \\
\bottomrule
\end{tabular}
\end{table}

Total state transfer overhead is $\sim$27\,ms for both models, confirming
that the binary serialization format is efficient and does not bottleneck
the hybrid pipeline.

\subsection{Batch vs.\ Sequential Prefill Comparison}

Table~\ref{tab:batch-vs-seq} compares the batch prefill results from this
work with the sequential dispatch measurements from our prior
study~\cite{atomgradient2026hybrid}.

\begin{table}[h]
\centering
\caption{Batch prefill (this work) vs.\ sequential dispatch
  (prior work~\cite{atomgradient2026hybrid}) for both models at 74 tokens
  on M2 Ultra.}
\label{tab:batch-vs-seq}
\begin{tabular}{llrrr}
\toprule
Model & Dispatch mode & Prefill (tok/s) & Speedup & TTFT \\
\midrule
\multirow{3}{*}{0.8B}
  & Sequential (prior) & 23.7 & 1.0$\times$ & $\sim$3{,}122\,ms \\
  & Batch (this work)  & 268.0 & \textbf{11.3$\times$} & $\sim$276\,ms \\
  & GPU baseline (MLX) & 736.2 & 31.1$\times$ & $\sim$101\,ms \\
\midrule
\multirow{3}{*}{2B}
  & Sequential (prior) & 23.7 & 1.0$\times$ & $\sim$3{,}122\,ms \\
  & Batch (this work)  & 172.7 & \textbf{7.3$\times$} & $\sim$429\,ms \\
  & GPU baseline (MLX) & 828.6 & 35.0$\times$ & $\sim$89\,ms \\
\bottomrule
\end{tabular}
\end{table}

For 0.8B, batch prefill achieves an 11.3$\times$ speedup over sequential
dispatch, reducing the gap with GPU prefill from 31.1$\times$ to
2.75$\times$ at 74 tokens. For 2B, batch prefill yields a 7.3$\times$
speedup, with GPU maintaining a larger lead (4.8$\times$ over batch)
due to the 2B model's higher compute requirements.

% ══════════════════════════════════════════════════════════════════════════════
\section{Discussion}
\label{sec:discussion}

\subsection{Prefill: ANE Batch vs.\ GPU}

ANE batch prefill \emph{outperforms} GPU prefill on short prompts: 149
vs.\ 138\,tok/s for 0.8B at 13 tokens. This advantage diminishes as
prompt length increases --- GPU prefill scales more efficiently with batch
size, reaching 736\,tok/s at 74 tokens while ANE saturates at 268\,tok/s.

The crossover point lies between 13 and 28 tokens for 0.8B. Below this
threshold, ANE batch prefill is strictly faster than GPU. For the 2B model,
GPU prefill already leads at 13 tokens (180 vs.\ 92\,tok/s), with the gap
widening at longer prompts.

This \textbf{short-prompt advantage} is significant for interactive
applications (chatbots, code completion) where prompts are typically
10--30 tokens.

\subsection{Decode Throughput Gap}

Pure ANE decode ($\sim$18\,tok/s for 0.8B, $\sim$12\,tok/s for 2B) is
significantly slower than GPU decode ($\sim$61--65\,tok/s for 0.8B,
$\sim$91--94\,tok/s for 2B). This confirms the finding from our prior
work that decode is memory-bandwidth-bound, and the GPU's 800\,GB/s
bandwidth advantage is decisive.

Hybrid decode (39--54\,tok/s for 0.8B) falls between pure ANE and pure
MLX. The gap from pure MLX is due to the state reconstruction overhead
in the MLX cache objects, not the serialization I/O (which is $<$30\,ms).

\subsection{The Concurrent Pipeline Argument}

The hybrid pipeline's total wall-clock time (6{,}050--6{,}335\,ms for
0.8B) is higher than pure MLX (3{,}679--3{,}846\,ms) when measured as a
single request. This is expected --- the sequential execution adds the
state transfer overhead.

However, the \textbf{real value} of ANE batch prefill is concurrent
execution. Since ANE prefill consumes only $\sim$0.22\,W of GPU power
(measured in~\cite{atomgradient2026hybrid}), the GPU is available for
\emph{simultaneously} decoding a previous request at full speed
($\sim$90\,tok/s). In a multi-request scenario:

\begin{itemize}[nosep]
  \item \textbf{Request $k$}: GPU decoding at $\sim$90\,tok/s
  \item \textbf{Request $k+1$}: ANE batch prefilling at 149--268\,tok/s
  \item \textbf{Net effect}: Zero idle time between requests; no GPU
        contention during prefill
\end{itemize}

This pipelining is impossible with GPU-only inference, where prefill and
decode must be serialized on the same hardware.

\subsection{Power and Thermal Implications}

Combining the power measurements from our prior work with the throughput
results here:

\begin{itemize}[nosep]
  \item \textbf{GPU prefill}: $\sim$62\,W GPU power (compute-bound)
  \item \textbf{ANE batch prefill}: $\sim$0.22\,W GPU + $\sim$1.58\,W ANE
        = $\sim$1.8\,W total
  \item \textbf{GPU decode}: $\sim$14\,W GPU power (bandwidth-bound)
\end{itemize}

During concurrent ANE prefill + GPU decode, total power is approximately
$1.8 + 14 = 15.8$\,W, compared to $62 + 14 = 76$\,W if both prefill and
decode used the GPU (hypothetical sequential scenario with two GPU
requests). This represents a \textbf{79\% power reduction} for the
prefill component.

On mobile devices (A-series, TDP $\sim$3--8\,W), GPU prefill at 62\,W
would immediately trigger thermal throttling. ANE prefill at $\sim$1.8\,W
fits within the mobile thermal envelope, enabling sustained inference
without clock-speed degradation.

\subsection{Limitations}

\begin{enumerate}[nosep]
  \item \textbf{Private API dependency}: The
        \texttt{AppleNeuralEngine.framework} API is undocumented and may
        change across macOS versions.
  \item \textbf{Hybrid decode gap}: Hybrid decode (39--54\,tok/s) does not
        match pure MLX decode ($\sim$61--94\,tok/s), likely due to cache
        reconstruction overhead in MLX's Python layer.
  \item \textbf{Single-machine results}: All measurements are from a
        single M2 Ultra; scaling behavior on other Apple Silicon variants
        (M1, M3, M4, A-series) remains to be validated.
  \item \textbf{No concurrent measurement}: We measured ANE prefill and
        GPU decode separately; true concurrent execution was not yet
        benchmarked end-to-end.
\end{enumerate}

% ══════════════════════════════════════════════════════════════════════════════
\section{Conclusion}
\label{sec:conclusion}

We demonstrated that ANE batch prefill via the private
\texttt{AppleNeuralEngine.framework} API achieves an 11.3$\times$
speedup over sequential dispatch, reaching 268\,tok/s for Qwen3.5-0.8B
on M2 Ultra. On short prompts ($\leq$13 tokens), ANE batch prefill
\emph{outperforms} GPU prefill (149 vs.\ 138\,tok/s).

Combined with the power measurements from our prior work (0.22\,W GPU
during ANE prefill vs.\ 62\,W during GPU prefill), these results validate
the feasibility of \textbf{concurrent ANE prefill + GPU decode} for
on-device multi-request serving. The state transfer overhead ($<$30\,ms)
is negligible, and the ANE's near-zero GPU power consumption during
prefill enables true hardware-level parallelism without thermal
contention.

Future work includes end-to-end concurrent pipeline benchmarking,
extension to larger models (9B+), and investigation of the ANE batch size
ceiling beyond \texttt{ANE\_SPATIAL=32}.

% ══════════════════════════════════════════════════════════════════════════════
\section*{Acknowledgments}

Experiments were conducted on a Mac Studio (M2 Ultra, 192\,GB). Power
measurements referenced from our prior study used \texttt{powermetrics}
on the same machine.

This work builds upon the foundational ANE reverse-engineering efforts
by \textbf{maderix}~\cite{maderix2024ane} and
\textbf{johnmai-dev}~\cite{johnmai2024anelm}. Without their pioneering
work on the private \texttt{AppleNeuralEngine.framework} API, this
research would not have been possible.

\begin{thebibliography}{9}

\bibitem{atomgradient2026hybrid}
AtomGradient.
\textit{Disaggregated LLM Inference on Apple Silicon: CoreML ANE Prefill
+ MLX GPU Decode}.
2026.
\url{https://github.com/AtomGradient/hybrid-ane-mlx-bench}

\bibitem{maderix2024ane}
maderix.
\textit{ANE: Apple Neural Engine reverse-engineering and experiments}.
\url{https://github.com/maderix/ANE}

\bibitem{johnmai2024anelm}
johnmai-dev.
\textit{ANE-LM: LLM inference on Apple Neural Engine}.
\url{https://github.com/johnmai-dev/ANE-LM}

\end{thebibliography}

\end{document}
